\documentclass[11pt]{article}
\usepackage{amsmath, amssymb}
\usepackage{geometry}
\geometry{margin=1in}
\usepackage{fancyhdr}
\usepackage{lastpage}
\usepackage{parskip} % Better paragraph spacing
\usepackage{titlesec} % Better section formatting
\usepackage{graphicx}
\usepackage{float}
\usepackage[most]{tcolorbox}
\usepackage{lmodern}
\usepackage{caption}
\usepackage{hyperref}
\usepackage{listings}
\usepackage{xcolor}
\usepackage{tikz}
\usetikzlibrary{automata, positioning, arrows.meta}

\definecolor{codegreen}{rgb}{0,0.6,0}
\definecolor{codegray}{rgb}{0.5,0.5,0.5}
\definecolor{codepurple}{rgb}{0.58,0,0.82}
\definecolor{backcolour}{rgb}{0.95,0.95,0.92}

% Define code styling
\lstdefinestyle{mystyle}{
    backgroundcolor=\color{backcolour},
    commentstyle=\color{codegreen},
    keywordstyle=\color{magenta},
    numberstyle=\tiny\color{codegray},
    stringstyle=\color{codepurple},
    basicstyle=\ttfamily\footnotesize,
    breakatwhitespace=false,
    breaklines=true,
    captionpos=b,
    keepspaces=true,
    numbers=left,
    numbersep=5pt,
    showspaces=false,
    showstringspaces=false,
    showtabs=false,
    tabsize=2
}

\lstset{style=mystyle}

% Define commands for reusable content
\newcommand{\coursename}{CS-550 Formal Verification}
\newcommand{\lecturenum}{07}
\newcommand{\lecturedate}{\today}
\newcommand{\professor}{Viktor Kuncak}

\hypersetup{colorlinks=true,linkcolor=blue,urlcolor=blue}

% Configure section spacing
\titlespacing*{\section}{0pt}{2ex plus 1ex minus .2ex}{1.5ex plus .2ex}
\titlespacing*{\subsection}{0pt}{1.5ex plus 1ex minus .2ex}{1ex plus .2ex}
% Header / Footer
\pagestyle{fancy}
\fancyhf{}
\fancyhead[L]{\coursename}
\fancyhead[R]{Lecture \lecturenum}
\fancyfoot[L]{EPFL - Computer Science}
\fancyfoot[C]{Author: Jonathan Dutoit}
\fancyfoot[R]{Page \thepage\ of \pageref{LastPage}}

\begin{document}

\begin{center}
  {\large EPFL - Computer Science} \\[0.5em]
  {\Large \textbf{\coursename}} \\[0.5em]
  {\large Lecture \lecturenum\ --- \lecturedate} \\[0.5em]
  {\large Author: Jonathan Dutoit} \\[0.5em]
  {\large Professor: \professor} \\[0.5em]
\end{center}

\vspace{1.5cm}

\begin{center}
  {\LARGE \textbf{Quantifier Elimination for Presburger Arithmetic}}
\end{center}

\vspace{1cm}

\section*{Formal Verification Methodology}

Formal verification of program as a three-step approach: 
\begin{enumerate}
	\item Express properties in logic or code \\ 
		(assertions, preconditions, post-conditions, invariants, run-time error conditions) 
	\item Compile program meaning to \textbf{logical formulas}  \\
	(verification-condition generator, symbolic execution)
\item Develop and use an automated theorem prover for generated conditions \\
	(SAT and SMT solving, resolution-based theorem proving, proof assistants)
\end{enumerate}

Which logic to use? 

\section*{Presburger arithmetic}

Integer arithmetic with logical operations (and, or, not), quantifiers, only addition as arithmetic operations, and $<$ and $=$ as a relation. 
\begin{itemize}
	\item minimalistically one can define a variant of it over non-negative natural numbers as having $\land, \neg, \forall, +, =$ as the only symbols. 
\end{itemize}

\begin{tcolorbox}[
    enhanced,
    frame hidden,
    borderline west={2pt}{0pt}{black},
    boxsep=5pt,
    left=8pt
]
\textbf{Definition.} One-Point Rule: If $\bar{y}$ is a tuple of variables not containing $x$, then $$\exists x. (x = t(\bar{y}) \land F(x, \bar{y})) \Leftrightarrow F(t(\bar{y}), \bar{y})$$
\begin{itemize}
	\item \textit{One point rule}: replaces LHS of equivalence by the RHS
	\item \textit{Flattening}: used when \textit{t} is complex, replaces RHS BY LHS.  \\

\end{itemize}
By the same reasoning, we have \textbf{dual one-point rule for} $\forall$ $$\neg (\exists x. (x = t(\bar{y}) \land F(x, \bar{y})) \Leftrightarrow F(t(\bar{y}), \bar{y}))$$ which is equivalent to $$\forall x. (x = t(\bar{y}) \to \neg F(x, \bar{y})) \Leftrightarrow \neg F(t(\bar{y}), \bar{y})$$ 

Hence we can deduce $$\forall x. (x = t(\bar{y}) \to F(x, \bar{y})) \Leftrightarrow F(t(\bar{y}), \bar{y})$$
\end{tcolorbox}

\subsection*{Quantifier Elimination (QE)}
Given a formula $F(\bar{y})$ containing quantifiers, find a formula $G(\bar{y})$ \textbf{equivalent} to $F(\bar{y})$, that has \textbf{no quantifiers}, and has a \textbf{subset (or equal set) of free variables} of \textit{F}.


\begin{itemize}
	\item To check satisfiability of $H(\bar{y})$: eliminate the quantifiers from $\exists \bar{y}. H(\bar{y})$ and evaluate. 
	\item Validity: eliminate quantifiers from $\forall \bar{y}. H(\bar{y})$ and evaluate. 
\end{itemize}

\begin{tcolorbox}[
  colback=gray!15,     % background color
  colframe=black!50,    % border color
  boxrule=1pt,       % border thickness
  arc=3mm,             % rounded corners
  left=4mm, right=4mm, top=2mm, bottom=2mm % padding
]
\textbf{Does Presburger Arithmetic admit QE?}
\vspace{0.3cm}

Depends on the particular set of symbols! Quantifier elimination implies: if you can define some relationship between variables using an arbitrary, possibly quantified formula \textit{F}, then you can also define same \textit{r} using another quantifier-free formula G. \\ 

\textbf{Example}: $\exists y. x = y + y + y$ only have an equivalent formula without quantifiers if we have \textit{divisibility}. 
\end{tcolorbox}

\subsection*{Presburger Arithmetic (PA)}
\[
	\{+, -, =, <\} \cup \{K | K \in \mathbb{Z}\} \cup \{(K \cdot \textbf{\_}) \, | \, K \in \mathbb{Z} \} \cup \{(K | \_ ) \, | \, K \in \mathbb{Z}\}
.\]

We also have, as usual: $\land, \lor, \neg, \to, \forall, \exists$.

\begin{itemize}
	\item \textbf{Atomic formula}: a relation applied to argument. \\
	$t_1 = t_2$, $t_1 < t_2$, $ K \, | \, t$
	\item \textbf{Literal}: Atomic formula or it's negation \\
	$\neg (x = y + 1)$
\end{itemize}

First, we consider the quantifier elimination problem of the form: $$\exists y. \quad L_1 \land ... \land L_n$$ This will prove to be sufficient to eliminate all quantifiers!

The question becomes: can we eliminate $\exists$ from any \textbf{quantifier-free formula}? $$\exists x. F(x, \bar{y})$$ where \textit{F} is quantifier free, meaning that it only has $\land, \lor, \neg$ applied to atomic formulas. Convert \textit{F} to \textbf{disjunctive normal form} where each $C_i$ is a \textbf{conjunction of literals}. $$\exists x. F \Leftrightarrow [\exists x. \bigvee^{m}_{i=1} C_i] \Leftrightarrow \bigvee^{m}_{i=1} (\exists x. C_i) $$

\begin{minipage}[t]{0.46\textwidth}
\subsubsection*{Nested Existential Quantifiers}
Eliminate one quantifier at a time:
\[
\begin{align}
\exists x_1.\exists x_2.\exists x_3.\,F_0(x_1,x_2,x_3,\bar{y}) \\
\rightsquigarrow \exists x_1.\exists x_2.\,F_1(x_1,x_2,\bar{y}) \\
\rightsquigarrow \exists x_1.\,F_2(x_1,\bar{y}) \\
\rightsquigarrow F_3(\bar{y})
\end{align}
\]
Each step eliminates $\exists$ from a conjunction of literals (after DNF).
\end{minipage}\hfill
\vrule
\hfill
\begin{minipage}[t]{0.46\textwidth}
\subsubsection*{Universal Quantifiers}
Reduce $\forall$ to the already handled existential case:
\[
\forall x.\,F_0(x,\bar{y})
\;\Leftrightarrow\;
\neg\big[\exists x.\,\neg F_0(x,\bar{y})\big]
\]
$\neg F_0$ is quantifier-free, so
$\exists x.\,\neg F_0(x,\bar{y}) \rightsquigarrow F_1(\bar{y})$, and therefore
\[
\forall x.\,F_0(x,\bar{y})
\;\rightsquigarrow\;
\neg F_1(\bar{y})
\]
\end{minipage}

\subsubsection*{Removing Any Alternation of Quantifiers}
Push negations in so that only $\exists$ blocks remain, alternating with negations:
\begin{align*}
&\exists x_1.\,\forall x_2.\,\forall x_3.\,\exists x_4.\;F_0(x_1,x_2,x_3,x_4,\bar{y})\\
\rightsquigarrow\;&\exists x_1.\,\neg\exists x_2.\,\exists x_3.\,\neg\exists x_4.\;F_0(x_1,x_2,x_3,x_4,\bar{y})\\
\rightsquigarrow\;&\exists x_1.\,\neg\exists x_2.\,\exists x_3.\;\neg F_1(x_1,x_2,x_3,\bar{y})\\
\rightsquigarrow\;&\exists x_1.\,\neg\exists x_2.\;F_2(x_1,x_2,\bar{y})\\
\rightsquigarrow\;&\exists x_1.\;\neg F_3(x_1,\bar{y})\\
\rightsquigarrow\;&F_4(\bar{y})
\end{align*}
Each quantifier alternation involves one disjunctive normal form transformation.
Note that for integers: $$\neg(x < y) \Leftrightarrow y \leq x \quad \quad x < y \Leftrightarrow x + 1 \leq y \quad \quad x \leq y \Leftrightarrow x < y + 1$$ Instead of $\leq$ we chose to use $<$ only. We do not write $x > y$ but only $y < x$, since y is the variable being quantified. 

\textbf{Normal form of terms}: All \textit{terms} are build from $K, +, -, K \cdot \_$, so using standard transformations they can be represented as: $K_0 + \sum^{n}_{i = 1} K_i x_i$. We call such term a \textit{linear term}.

\textbf{Normal form for literals in PA}:

\[
\begin{array}{l}
	\neg (t_1 < t_2) \quad \text{becomes} \quad t_2 < t_1 + 1 \\
	\neg (t_1 = t_2) \quad \text{becomes} \quad t_1 < t_2 \lor t_2 < t_1 \\ 
	t_1 = t_2 \quad \text{becomes} \quad t_1 < t_2 + 1 \land t_2 < t_1 + 1 \quad (*) \\ 
	\neg (K \, | \, t) \quad \text{becomes} \quad \bigvee^{K-1}_{i=1} K \, | \, t + i \\
	t_1 < t_2 \quad \text{becomes} \quad 0 < t_2 - t_1
\end{array}
\] 
To remove generated disjunctions, compute DNF again. \\
(*) Transformed equalities for simplicity. Usually, handle them directly. 

\textbf{Example:} eliminate \textit{y} from this formula $$\exists y. \, 3y - 2w + 1 > -w \land 2y - 6 < z \land 4 | 5y + 1$$

Simplify/normalize using properties of integer operations: $$\exists y. \quad 0 < -w + \textbf{3y} + 1 \quad \land \quad 0 < \textbf{-2y} + z + 6 \quad \land \quad 4 \, | \, \textbf{5y} + 1$$ Least common multiple of coefficients \textit{M = lcm(3, 2, 5) = 30} $$\exists y. \quad 0 < -10w + \underline{\textbf{30y}} + 10 \quad \land \quad 0 < \underline{\textbf{-30y}} + 15z + 90 \quad \land \quad 24 \, | \, \underline{\textbf{30y}} + 6$$ 

Let \textit{x} denotes \textit{30y}. It is \textbf{not an arbitrary x}. It is divisible by 30. $$ \exists x. \quad 0 < -10w + x + 10 \, \land \, 0 < -x + 15z + 90 \, \land \, 24 \, | \, x + 6 \, \land \, 30 \, | \, x$$  

\begin{tcolorbox}[
  colback=blue!15,     % background color
  colframe=black!50,    % border color
  boxrule=1pt,       % border thickness
  arc=3mm,             % rounded corners
  left=4mm, right=4mm, top=2mm, bottom=2mm % padding
]
\textbf{Lower and upper bounds}
\vspace{0.3cm}

Consider the coefficient next to \textit{x} in $0 < t$. If it is -1, move the term to left side. If it is 1, move the remaining terms to the left side. We obtain formula $F_1(x)$ of the form: $$\bigwedge^{L}_{i=1}a_i < x \land \bigwedge^{U}_{j=1} x < b_j \land \bigwedge^{D}_{i=1} K_i \, | \, (x + t_i)$$ If there are no divisibility constraints (\textit{D = 0}), the formula is equivalent to: $$\mathop{\mathrm{max}}\limits_{i} a_i + 1 \leq \mathop{\mathrm{max}}\limits_{j}b_j - 1 \Leftrightarrow \bigwedge_{ij} a_i + 1 < b_j$$
\end{tcolorbox}

\subsection*{Replacing variable by test terms}

\textit{\underline{Common technique in quantifier elimination}}: replacing $F_1(x)$ with $\bigvee_k F_1(t_k)$, where $t_k$ do not contain \textit{x}. If $F_1(t_k)$ holds then certainly $\exists x. F_1(x)$. The example terms $t_i$ when $D = 0$ and $L > 0$ $$\bigvee_{k=1}^{L} F_1(a_k + 1) \quad \quad \text{ensuring that at least one of them evalutes to max} a_i + 1$$ and if $D>0$, i.e., we have additional divisibility constraints: $$\bigvee^{L}_{k=1} \bigvee^{N}_{i=1} F_1(a_k + i) \quad \quad \text{where N is the lcm of }  K_1, ..., K_D$$

$$ \exists x. F_1(x) \quad \equiv \quad \exists x. \, -10 + 10w < x \, \land \, x < 15z + 90 \, \land \, 24 \, | \, x + 6 \, \land \, 30 \, | \, x$$

Let's denote the lower bound $a_1 = -10w + 10$, $b_1 = 15z + 90$, $K_1 = 24, K_2 = 30$ hence lcm($K_1, K_2$) = 120). 

We replace with $a_1 + i < b_1$, meaning that x is replaced with $-10 + 10w +i$, which corresponds to: 
\[
	\bigvee^{120}_{i=1} 10w + i < 100 + 15z \, \land \, 24 \, | \, 10w - 4 + i \, \land \, 30 \, | \, 10w - 10 + i
\]
This is the solution. $i$ is convenient written so that we avoid writing 120 disjuncts, but it's not a \textit{variable}. 

\subsection*{No lower bounds}

Consider the case $L = 0$. $$ \bigwedge^{U}_{j=1} x < b_j \land \bigwedge^{D}_{i=1} K_i \, | \, (x + t_i)$$ We cannot try all lower bounds, as there are none. Thus, we can always find a $x$ small enough to satisfy the upper bounds. Hence, we drop all constraints except divisibility, obtaining $F_2(x)$: $$\bigwedge^{D}_{i=1} K_i \, | \, (x + t_i)$$ and then eliminate quantifier as $$\bigvee^{N}_{i=1} F_2(i)$$. \textbf{Example} $$\exists x. \, x < 5 + 7z \, \land \, 2 \, | \, x \, \land 3 \, | \, x + 2$$ N = lcm(2, 3) = 6. $$\bigvee^{6}_{i=1} (2 \, | \, i \, \land \,  3 \, | \, i + 2)$$

\section*{Summary}

This complete the description of a quantifier elimination algorithm for Presburger Arithmetic (PA). 

This algorithm and its correctness prove that: 
\begin{itemize}
	\item PA admits quantifier elimination
	\item Satisfiability, validity, entailment, equivalence of PA formulas is decidable. \\
	We can use the algorithm to prove verification conditions. \\
	Not the most efficient way, but gives insights on what can later be build to come up with better algorithms
	\item Quantified and quantifier-free formulas have the same expressive power. 
\end{itemize}

Regardless of the procedure, cannot escape inherent computational complexity of deciding if quantified Presburger arithmetic formula is true, which has been characterized using alternating complexity classes by Bernman. 

A set \textit{A} is in \textit{STA(s(n), t(n), a(n))} iff there exists a single-tape \textit{alternating} Turing machine $M_i$ (making $\land$ as well as $\lor$ choices) that accepts \textit{A} and runs withing space \textit{s(n)}, time \textit{t(n)} and performs at most \textit{a(n)} alternations between $\land$ and $\lor$ choices. 

Deciding of a PA formula of size \textit{n} with \textit{m} quantifiers is in $$STA(*, 2^{n^{O(m)}}, m)$$ Complexity is particularly sensitive to \textit{m}. 

\end{document}
